\section{メンテナンスとトラブルシューティング}

\subsection{フィルターの交換}
使用している $1.0\,\mu\mathrm{m}$ Mylarフィルターは極めて薄く、微小な圧力差でも破れる。
\begin{itemize}
    \item 真空開放時: 必ずスローリークバルブを使用し、数分かけて大気圧に戻すこと。急激な開放は厳禁。
    \item 交換時期: 破れが目視で確認された場合、またはMylar画像の信号強度が異常に上昇した場合（可視光漏れ）。
\end{itemize}

\subsection{MCPの保管}
MCPは湿気に弱いため、長期間使用しない場合は以下のいずれかの処置をとる。
\begin{enumerate}
    \item 真空引きを継続する（推奨）。
    \item 取り外してデシケーター（防湿庫）内で保管する。
\end{enumerate}